Pendiente de redacción completa con bibliografía verificada. Este apartado debe cubrir modelado de UAV multirrotor, control en cascada PID, simulación dinámica de cuadricópteros, aprendizaje por imitación y arquitecturas MLP, GRU y LSTM aplicadas a control.

Como base inicial, el modelado y control de aeronaves no tripuladas se apoyará en textos de referencia de sistemas UAV~\cite{beard2012small}. Para cuadricópteros, se revisarán formulaciones de generación y seguimiento de trayectorias como la propuesta por Mellinger y Kumar~\cite{mellinger2011minimum}. La parte de aprendizaje por imitación se conectará con métodos supervisados y de agregación de datos como DAgger~\cite{ross2011dagger}, dejando claro que el alcance del TFG es imitación supervisada y no aprendizaje por refuerzo.

La comparación de simuladores debe justificar la decisión de implementar un simulador propio sin convertir el apartado en una revisión de herramientas de software. La discusión debe centrarse en trazabilidad física, control de hipótesis y adecuación al objetivo del TFG.
